% ===== UNIT 13: Stream 1: augmentation + hyper-parameters (3.2.3-3.2.4) =====
% Status: AI-DRAFT  (change to REWRITTEN / REVIEWED when done; see CHECKLIST.md)
% Rules: keep every \cite, \ref, \label, number and table value unless you checked it against the logs.
\subsubsection{Augmentation pipeline}
Table~\ref{tab:augmentation} lists the full augmentation sequence applied
during Stream~1 training. All reported Stream~1 results use a
224\,$\times$\,224 input; hue/colour-corrupting augmentations (random
grayscale, strong hue jitter) were excluded as they remove vehicle colour
identity cues. A controlled 256\,$\times$\,256 \emph{training}-resolution
comparison was not reproduced in this campaign; an \emph{inference}-resolution
sensitivity analysis (Section~\ref{sec:res_sens}) shows the deployed checkpoint
performs best at its native 224\,$\times$\,224.

\begin{table}[!t]
  \centering
  \caption{Stream~1 training augmentation pipeline. Steps are applied in the
  listed order to every training image}
  \label{tab:augmentation}
  \small
  \begin{tabular}{clL{8.4cm}}
    \toprule
    Step & Transform & Parameters \\
    \midrule
    1 & Resize              & $224\times224$, bicubic \\
    2 & Random H-flip       & $p=0.5$ \\
    3 & Pad $+$ rand.\ crop & padding $=10$\,px \\
    4 & ColorJitter         & brightness $=0.2$, contrast $=0.15$,
                              saturation $=0.1$, hue $=0.0$ \\
    5 & Random Gauss.\ blur & kernel $=5$, $p=0.2$ \\
    6 & Random perspective  & distortion $=0.1$, $p=0.2$ \\
    7 & Normalise           & CLIP $\mu$=(0.481,\,0.458,\,0.408),
                              $\sigma$=(0.269,\,0.261,\,0.276)~\cite{radford2021clip} \\
    8 & Random erasing      & $p=0.5$, scale $=[0.02,\,0.33]$ \\
    \bottomrule
  \end{tabular}
\end{table}

\subsubsection{Training hyperparameters}
Table~\ref{tab:hyperparams} compares the proposed Stream~1 configuration with
the published TransReID recipe~\cite{he2021transreid}. Training required
approximately 2.5--3.0 GPU-hours on $2\times$ NVIDIA T4 GPUs with DataParallel.
The exact $P{\times}K$ tuple is not recorded in the training logs; only the
effective batch size of 96 is confirmed. For future retraining we recommend
$P{\times}K = 12{\times}8$ as the default consistent with this batch
size---stated explicitly as a recommendation rather than a recovered record;
retraining reproducibility under the released configuration is bounded by the
$\pm 0.55$\,pp seed band of Section~\ref{sec:seed_variance}.

\begin{table}[!t]
  \centering
  \caption{Training hyperparameters: Stream~1 versus
  TransReID~\cite{he2021transreid}. LLRD denotes layer-wise learning-rate
  decay}
  \label{tab:hyperparams}
  \small
  \begin{tabular}{L{4.0cm}L{4.4cm}L{4.6cm}}
    \toprule
    Hyperparameter & TransReID~\cite{he2021transreid} & Ours (Stream~1) \\
    \midrule
    Backbone     & ViT-B/16 $+$ SIE $+$ JPM       & ViT-B/16 $+$ SIE $+$ JPM \\
    Pre-training & ImageNet-21k$\!\to\!$1K        & OpenAI CLIP~\cite{radford2021clip} \\
    Input size   & $256\times256$                 & $224\times224$ \\
    Optimiser    & SGD (momentum 0.9)             & AdamW~\cite{loshchilov2019adamw} \\
    LR backbone  & 0.008                          & $3.5\times10^{-4}$ (LLRD $=0.75$) \\
    LR head      & 0.008                          & $3.5\times10^{-3}$ \\
    LR schedule  & Cosine decay                   & Cosine $+$ 10-epoch warm-up \\
    Batch size   & 64 ($16\times4$)               & 96 ($2\times$T4 DataParallel) \\
    Epochs       & 120                            & 140 \\
    Loss         & CE $+$ triplet                 & CE-LS $+$ triplet $+$ centre \\
    Mixed prec.  & fp16                           & fp16 \\
    SIE cameras  & 20                             & 20 \\
    JPM groups   & 4                              & 4 \\
    BNNeck       & Yes                            & Yes \\
    \bottomrule
  \end{tabular}
\end{table}
